%! TeX program = lualatex
\documentclass[tikz, crop]{standalone}
\input{../typography.tex.preamble}
\input{../tikz.tex.preamble}

\pgfplotsset{
  extrema/.style = {
    basic curve,
    width={200pt},
    xlabel={\empty},
    ylabel={\empty},
    enlargelimits,
    xmin=-1, xmax=3, ymin=-1, ymax=3,
  }
}

\begin{document}


\begin{tikzpicture}
  \begin{axis}[extrema]
    \addplot[very thick, domain=-1:3, smooth, samples=1000] {abs(sin(pi*(x-1))*(x-1))};
    \draw[fill=black] (axis cs:-1,0) circle[radius=0.1];
    \draw[fill=black] (axis cs: 3,0) circle[radius=0.1];
  \end{axis}
\end{tikzpicture}

\begin{tikzpicture}
  \begin{axis}[extrema]
    \addplot[very thick] coordinates {(0,1) (1,3)};
    \addplot[very thick] coordinates {(1,1) (2,0)};
    \draw[fill=black] (axis cs:0,1) circle[radius=0.1];
    \draw[fill=white] (axis cs:1,3) circle[radius=0.1];
    \draw[fill=black] (axis cs:1,1) circle[radius=0.1];
    \draw[fill=black] (axis cs:2,0) circle[radius=0.1];
  \end{axis}
\end{tikzpicture}

\begin{tikzpicture}
  \begin{axis}[extrema]
    % shifted half a unit right, so this domain starts at 1/2 like f_3 and f_5
    \addplot[very thick] coordinates {(1/2,1) (3/2,2)};
    \addplot[very thick] coordinates {(3/2,1) (5/2,0)};
    \draw[fill=black] (axis cs:1/2,1) circle[radius=0.1];
    \draw[fill=white] (axis cs:3/2,2) circle[radius=0.1];
    \draw[fill=black] (axis cs:3/2,1) circle[radius=0.1];
    \draw[fill=white] (axis cs:5/2,0) circle[radius=0.1];
  \end{axis}
\end{tikzpicture}

% Defined on [-1, infinity), oscillating about the horizontal asymptote
% y = 3/2 with an amplitude that decays, so the values settle towards 3/2
% without the function ever being equal to it.  The decay rate is 0.6 rather
% than 1.5: over a four-unit window the faster rate flattens the curve onto the
% asymptote within the first unit and the rest of the sketch says nothing.  The
% plot runs past the right edge of the axis so the frame clips it, which reads
% as the domain continuing rather than the function stopping at 3.
\begin{tikzpicture}
  \begin{axis}[extrema]
    % the asymptote goes down first so the curve draws over it
    \draw[dashed] (axis cs:-1.5,3/2) -- (axis cs:3.5,3/2);
    \addplot[very thick, smooth, samples=400, domain=-1:3.5]
    {e^(-0.6*(x+1))*cos(2*pi*(x+1)) + 3/2};
    \draw[fill=black] (axis cs:-1,5/2) circle[radius=0.1];
  \end{axis}
\end{tikzpicture}

% Vertical asymptotes at both open ends, so again neither extremum exists.
% Shifted half a unit right, matching f_3, so the left asymptote sits at 1/2.
% The plotted domain stops where the curve is already past the top and bottom
% of the axis, so both asymptotes are clipped by the frame.
\begin{tikzpicture}
  \begin{axis}[extrema]
    \draw[dashed] (axis cs:1/2,-1.5) -- (axis cs:1/2,3.5);
    \draw[dashed] (axis cs:5/2,-1.5) -- (axis cs:5/2,3.5);
    % the denominator's factor of 2 is dropped for the same reason as in f_3:
    % it doubles the distance between the branches and the asymptotes
    \addplot[very thick, smooth, samples=1000, domain=0.62:2.38]
    {(x-3/2)/((x-1/2)*(5/2-x)) + 1};
  \end{axis}
\end{tikzpicture}

\begin{tikzpicture}
  \begin{axis}[extrema]
    % vertical asymptote at 1/2, horizontal asymptote at 2
    \draw[dashed] (axis cs:1/2,-1.5) -- (axis cs:1/2,3.5);
    \draw[dashed] (axis cs:-1.5,2) -- (axis cs:3.5,2);
    % 2 - 1/(x-1/2) rises towards 2 from below and is never equal to it, which
    % is the contrast with f_4: that one crosses its asymptote on every swing.
    % The coefficient 1 keeps the branch a visible distance from the vertical
    % asymptote, entering the frame about 0.3 units to its right, and the plot
    % runs past the right edge so the frame clips it rather than the curve
    % appearing to stop.
    \addplot[very thick, smooth, samples=400, domain=0.72:3.5]
    {2 - 1/(x-1/2)};
  \end{axis}
\end{tikzpicture}


\begin{tikzpicture}
  \begin{axis}[extrema, xmin=-2, xmax=4, ymin=-2, ymax=4, width={180pt}, xticklabels={\empty}, yticklabels={\empty}, enlargelimits=false]
  \end{axis}
\end{tikzpicture}
\end{document}
